% The appendix contains reproducibility detail and secondary analyses only.
% Claim decisions and their decisive bounds remain in the main paper.

\section{Reproducibility and Statistical Details}
\label{app:evaluation-contract}
\label{app:datasets}
\label{app:objectives}

The released campaign manifest is the source of truth for every request,
candidate, semantic group, checkpoint, seed, random pool, and configuration
hash. Table~\ref{tab:datasets} is a readable digest, not a replacement for
that manifest. The denominator is the predeclared matrix: planned rows remain
visible when they are never attempted, terminate early, fail integrity, do not
reach the common forgetting criterion, or lack complete all-arm support.

\begin{table*}[t]
\caption{\textbf{Compact campaign and outcome contract.} Panel A records the
predeclared scope. The three development folds are mutually disjoint and
target-free. Panel B defines constraint-matched forgetting and the common
utility floor; exact metrics and thresholds are frozen in the manifest before
target outcomes.}
\label{tab:datasets}
\label{tab:objective-config}
\label{tab:constraint-contract}
\centering
\scriptsize
\setlength{\tabcolsep}{2.7pt}
\resizebox{\textwidth}{!}{%
\begin{tabular}{@{}llllcccc@{}}
\toprule
\multicolumn{8}{@{}l}{\textbf{A. Predeclared setting scope}} \\
Dataset / role & Model / precision & $\mathcal D_{\rm cal}/\mathcal D_{\rm pred}/\mathcal D_{\rm prot}$ &
Target requests & Primary parents & Seeds & $|\mathcal C_{\rm disc}|/|\mathcal C_{\rm audit}|$ & Planned cells \\
\midrule
TOFU / primary & Qwen2.5-1.5B / fp32 & \tblph & \tblph & 7 & \tblph & \tblph & \tblph \\
TOFU / scale & Qwen2.5-7B / fp32 & \tblph & \tblph & 7 & \tblph & \tblph & \tblph \\
TOFU / family & Llama-3.1-8B / fp32 & \tblph & \tblph & 7 & \tblph & \tblph & \tblph \\
WMDP-bio / MMLU & \tblph & \tblph & \tblph & 7 & \tblph & \tblph & \tblph \\
MUSE-News & \tblph & \tblph & \tblph & 7 & \tblph & \tblph & \tblph \\
RWKU & \tblph & \tblph & \tblph & 7 & \tblph & \tblph & \tblph \\
MUSE-Books / stress & \tblph & \tblph & \tblph & 7 & \tblph & \tblph & \tblph \\
PISTOL / stress & \tblph & \tblph & \tblph & 7 & \tblph & \tblph & \tblph \\
\bottomrule
\end{tabular}}

\vspace{3pt}
\resizebox{0.88\textwidth}{!}{%
\begin{tabular}{@{}llllll@{}}
\toprule
\multicolumn{6}{@{}l}{\textbf{B. Shared final-checkpoint feasibility region}} \\
Audit & Metric & Favorable direction & Frozen boundary & Stopping role & Reported slack \\
\midrule
Direct forgetting & \tblph & \tblph & \tblph & parent entry + final constraint & yes \\
Paraphrase forgetting & \tblph & \tblph & \tblph & final constraint & yes \\
Extraction / generation & greedy autoregressive gold-token recall & lower & \tblph & final constraint & yes \\
General utility & \tblph & \tblph & \tblph & final floor & yes \\
Native retain audit & mean/CVaR$_{.95}$ $\Delta$NLL + dataset-native metric & lower/higher & frozen NI margin & outcome only & n.a. \\
\bottomrule
\end{tabular}}
\end{table*}

The primary objective roster is GradDiff, NPO, SimNPO, GRU, RMU, RepNoise,
and Circuit Breakers; GA and IdkDPO are stress controls. For each parent the
artifact records the implementation identifier, source commit, trainable
block, learning rate, update and processed-token budgets, saved-step grid,
forget/reference coefficients, stopping rule, and hashes of all frozen
inputs. Readout labels are descriptive and never select a mixture endpoint.

The candidate manifest freezes $\mathcal C$, the discovery/audit split by
semantic group, repair eligibility, and a score-independent random stream. Direct
answer duplicates, forget-fact paraphrases, and template-derived restatements
are ineligible for repair; entity overlap is retained as a disclosed stratum.
Only $\mathcal C_{\rm disc}$ fits empirical midrank maps or supplies repair
examples. Audit scores are computed at $\theta_0$ and sealed. Candidate IDs
break ties only when selecting the last member of a top-$K_p$ pool.
The neutral stream and differentiable guard sets are selector-independent.
General-utility checkpoint selection uses a separate frozen set; neither
$\mathcal C_{\rm audit}$ nor its native retained outcome is used by a guard,
stopping rule, or checkpoint selector.

\subsection{Units, funnels, and decisions}
\label{app:outcomes}

Every prediction correlation is formed within one
model--request--seed--parent cell. Seeds are averaged within requests and
requests receive equal weight. The paired hierarchical bootstrap resamples
requests, trajectory seeds within requests, and semantic groups within
requests. Its replicate count, seed, and
finite-sample percentile rule are manifest fields. No candidate is pooled
across requests before a correlation is computed.

Within a request, empirical $\mathrm{CVaR}_{.95}$ is the mean of the largest
5\% damage values, using fractional weight at the quantile boundary. When an
audit has fewer than 20 candidates, this reduces to the maximum and the row is
flagged as low-tail-support rather than pooled with larger audits.

The audit exposes four denominators. $N_{\rm prof}$ counts all predeclared
profiles and $N_{\rm valid}$ those passing the outcome-blind integrity check.
$N_{\rm att}$ counts all predeclared parent trajectories and $N_{\rm reach}$
those reaching the common entry criterion. $N_{\rm rv}$ is the
reached--valid intersection. Prediction common support requires joint score,
both endpoints, the frozen simple control, and damage on the same units. A
tail cell is eligible only when its positive-damage count is at least $m=K_p$,
with $q=m/|\mathcal C_{\rm audit}|$. Tail-ineligible cells remain in
the RQ1 correlation denominator; $N_{\rm tail}/N_{\rm rv}$ is always reported,
and the tail claim requires at least 0.80 coverage. Protection feasibility
requires mixture, no repair, every frozen random draw, $S_0$, and $S_1$ to
satisfy all four constraints on the same units; common support additionally
requires both damage outcomes and the native retained metric for every arm.
No missing outcome is imputed.

RQ1 is the intersection--union test (IUT)
$\rho_S>0\land g_G>0\land g_H>0\land L_{\rm tail}>0$, with the tail term
additionally requiring at least 0.80 coverage. RQ2 requires
$f_\rho>0.80$, $f_K>0.70$, $g_H>0$, and $g_{\rm ctl}>0$ through their
one-sided lower bounds. Protection is the IUT over four comparators, two
$\Delta$NLL outcomes, and one native-metric non-inferiority margin per
comparator. Thus success requires four RQ1 lower bounds, four RQ2 lower
bounds, and eight superiority plus four non-inferiority bounds for RQ3 to
cross their frozen boundaries in the favorable direction. Statistical
success and eligibility are separate
fields. A fallback weight, invalid profile, missing
expected cell, non-reach, infeasible arm, incomplete random-draw roster, or
incomplete common support forces claim pass to ``no,'' even if a conditional
effect is printable. Exact-gradient allocation is a same-estimand numerical
reference outside the protection IUT.

Parent rows use their one-sided IUT level without a within-row multiplicity
penalty. The later setting-level rule can select a passing parent from each
readout group, so that selection applies a Bonferroni threshold across the four
output-readout or three representation-readout parents. This correction is
used only for the breadth statement; it does not hide the complete row roster.

The paper-level breadth rule is defined in Section~\ref{sec:exp-supporting}; it
does not pool parent rows or turn the least-favorable row in
Table~\ref{tab:robustness} into a simultaneous confidence bound.

\FloatBarrier

\section{Additional Prediction Analyses}
\label{app:prediction-supp}

The main paper defines the tail premise and carries every parent-level claim
decision. Table~\ref{tab:tail-structure} provides the parent breakdown and
secondary controls. Tail support is conditioned on a realized parent damage
vector, not on profile validity, so a failed profiler cannot erase evidence
about whether aggregate retention hides a coherent tail. The permutation
keeps request-specific group sizes fixed. If a request has zero positive
damage mass, its concentration ratio is undefined: we report it as a
zero-mass request and exclude it only from that ratio, never from the
prediction denominator. Tail size and minimum support are frozen in the
manifest.

\begin{table*}[t]
\caption{\textbf{Additional tail, reference, and construction analyses.}
Panel A expands the premise behind Figure~\ref{fig:channel-main}A. Panel B
compares pre-unlearning references on equal parent/request support. Panel C
contains representative one-factor checks; the full grid and bootstrap draws
are artifacts. None of these panels replaces the main table's absolute and
two-endpoint prediction test.}
\label{tab:tail-structure}
\label{tab:prospective-prediction}
\label{tab:profile-construction}
\label{tab:horizon-sensitivity}
\centering
\scriptsize
\setlength{\tabcolsep}{2pt}
\resizebox{0.88\textwidth}{!}{%
\begin{tabular}{@{}lccccccc@{}}
\toprule
\multicolumn{8}{@{}l}{\textbf{A. Damage-tail evidence on repair-held-out audit behaviors}} \\
Parent & Mean damage & CVaR$_{.95}$ & $M_q/q$ [95\% CI] &
$L_{\rm grp}$ [95\% CI] & Permutation $p$ & Damage support $n/N$ & Tail $n$ \\
\midrule
GradDiff & \tblph & \tblph & \tblph & \tblph & \tblph & \tblph & \tblph \\
NPO & \tblph & \tblph & \tblph & \tblph & \tblph & \tblph & \tblph \\
SimNPO & \tblph & \tblph & \tblph & \tblph & \tblph & \tblph & \tblph \\
GRU & \tblph & \tblph & \tblph & \tblph & \tblph & \tblph & \tblph \\
RMU & \tblph & \tblph & \tblph & \tblph & \tblph & \tblph & \tblph \\
RepNoise & \tblph & \tblph & \tblph & \tblph & \tblph & \tblph & \tblph \\
CB & \tblph & \tblph & \tblph & \tblph & \tblph & \tblph & \tblph \\
\bottomrule
\end{tabular}}

\vspace{3pt}
\resizebox{\textwidth}{!}{%
\begin{tabular}{@{}llccccccc@{}}
\toprule
\multicolumn{9}{@{}l}{\textbf{B. Pre-unlearning reference roster; descriptive equal-parent summaries}} \\
Family & Predictor & Candidate reverse & All $\rho$ & Output $\rho$ &
Representation $\rho$ & Top-$q$ recall & Common support & Role \\
\midrule
Gradient & exact gradient energy & yes & \tblph & \tblph & \tblph & \tblph & \tblph & numerical reference \\
Directional & forget-direction symmetric FD & no & \tblph & \tblph & \tblph & \tblph & \tblph & baseline \\
Loss response & one random direction & no & \tblph & \tblph & \tblph & \tblph & \tblph & baseline \\
Representation & sentence-embedding proximity & no & \tblph & \tblph & \tblph & \tblph & \tblph & baseline \\
Lexical & token-overlap proximity & no & \tblph & \tblph & \tblph & \tblph & \tblph & baseline \\
Representation & request-shuffled proximity & no & \tblph & \tblph & \tblph & \tblph & \tblph & negative control \\
Loss & initial NLL & no & \tblph & \tblph & \tblph & \tblph & \tblph & baseline \\
Length & answer length & no & \tblph & \tblph & \tblph & \tblph & \tblph & baseline \\
Random & repeated frozen rankings & no & \tblph & \tblph & \tblph & \tblph & \tblph & chance \\
\bottomrule
\end{tabular}}

\vspace{3pt}
\resizebox{\textwidth}{!}{%
\begin{tabular}{@{}lllccccc@{}}
\toprule
\multicolumn{8}{@{}l}{\textbf{C. Representative one-factor checks}} \\
Axis & Primary / variant & Scope & Joint $\rho$ & min gain [min LB] &
Top-$q$ recall & Common support & Eligible/pass \\
\midrule
Representation layer & final / penultimate & one output + one representation parent & \tblph & \tblph & \tblph & \tblph & \tblph \\
Pooling & answer-token mean / final token & same & \tblph & \tblph & \tblph & \tblph & \tblph \\
Loss-shake block & frozen late MLP / last layer & same & \tblph & \tblph & \tblph & \tblph & \tblph \\
Mixture & parent-specific / one global weight & all primary parents & \tblph & \tblph & \tblph & \tblph & \tblph \\
Horizon & first reach / fixed or terminal & all reached parents & \tblph & \tblph & \tblph & \tblph & \tblph \\
\bottomrule
\end{tabular}}
\end{table*}

All reference predictors use the same candidate folds and outcome horizons.
Construction ablations keep the primary prediction weight fixed; a separately
reported development refit diagnoses weight mismatch but cannot replace that
ablation. The request-shuffled control draws a different forget request of the
same declared size before target outcomes.

\FloatBarrier

\section{Loss-Shake Numerical Validation}
\label{app:estimator-stability}
\label{app:probe-details}
\label{app:randomized-impl}

For each Gaussian draw, the implementation normalizes
$v_r=g_r/\lVert g_r\rVert_2$, saves the declared parameter block, applies
$\theta_0+\eta v_r$ and $\theta_0-\eta v_r$ from that same copy, and restores
it before the next direction. Candidate order,
batching, masks, padding, precision, and dropout state are identical within a
pair. Signed responses are accumulated in fp32 and persisted before squaring,
so nested-$R$, direction-bank, and split-half analyses require no new model
evaluation. The dimension factor in Eq.~\eqref{eq:loss-shake-identity} is
stored even though it does not affect ranking.

The confirmatory operating point uses fp32 parameters. Bfloat16 is retained
only as a numerical boundary check: a configuration is invalid when the
realized block displacement or changed-coordinate fraction falls below its
frozen perturbation-survival threshold.

The exact reference streams one candidate at a time, records its gradient
energy on the same block, and discards the vector. Outcome-blind development
calibration chooses the smallest $(R,\eta)$ satisfying the frozen criteria for
exact-rank fidelity, top-$K_p$ overlap, perturbation survival, and
direction-split stability. A
target profile is valid only when the latter forward-only checks pass on
$\mathcal C_{\rm disc}$ before the sealed audit is opened.
\label{app:fidelity-gate}

\begin{table*}[t]
\caption{\textbf{Loss-shake fidelity--cost frontier and numerical
sensitivity.} Timing uses device synchronization and a common measurement
boundary; memory is peak allocated device memory. The selected operating
point is marked frozen. Exact gradients are evaluated only as a reference.}
\label{tab:bwfree}
\label{tab:estimator-validation}
\centering
\scriptsize
\setlength{\tabcolsep}{2.2pt}
\resizebox{\textwidth}{!}{%
\begin{tabular}{@{}lllcccccc@{}}
\toprule
\multicolumn{9}{@{}l}{\textbf{A. Matched fidelity--cost frontier}} \\
Model / block / precision & Profiler & $R$ & $\rho$ vs. exact & overlap@$K_p$ &
Pool passes / candidate reverse & Time & Peak memory & Integrity valid $n/N$ \\
\midrule
Qwen2.5-1.5B / \tblph / fp32 & exact energy & --- & 1.00 & 1.00 & candidate F+B / yes & \tblph & \tblph & n.a. \\
& exact energy (batched, if feasible) & --- & 1.00 & 1.00 & batched F+B / yes & \tblph & \tblph & n.a. \\
& loss-shake & \tblph & \tblph & \tblph & $2R$ pool F / no & \tblph & \tblph & \tblph \\
& loss-shake (frozen) & \tblph & \tblph & \tblph & $2R$ pool F / no & \tblph & \tblph & \tblph \\
Qwen2.5-7B / \tblph / fp32 & exact energy & --- & 1.00 & 1.00 & candidate F+B / yes & \tblph & \tblph & n.a. \\
& loss-shake (frozen) & \tblph & \tblph & \tblph & $2R$ pool F / no & \tblph & \tblph & \tblph \\
Llama-3.1-8B / \tblph / fp32 & loss-shake (frozen) & \tblph & \tblph & \tblph & $2R$ pool F / no & \tblph & \tblph & \tblph \\
\bottomrule
\end{tabular}}

\vspace{3pt}
\resizebox{0.9\textwidth}{!}{%
\begin{tabular}{@{}llccccc@{}}
\toprule
\multicolumn{7}{@{}l}{\textbf{B. One-factor numerical sensitivity}} \\
Axis & Frozen / variants & $\rho$ range & overlap range & split-half range &
Valid $n/N$ & Boundary diagnosis \\
\midrule
Directions & $R$ / nested prefixes and new banks & \tblph & \tblph & \tblph & \tblph & \tblph \\
Radius & $\eta$ / log-grid neighbors & \tblph & \tblph & \tblph & \tblph & \tblph \\
Precision & declared / fp32, bf16 & \tblph & \tblph & \tblph & \tblph & \tblph \\
Pool size & full / nested pools, fixed common subset & \tblph & \tblph & \tblph & \tblph & \tblph \\
\bottomrule
\end{tabular}}
\end{table*}

Rank agreement validates computation of $\widehat q_G$ only. Fidelity alone
establishes neither component's predictive added value; the main paired gains
and simple-control comparison decide that separate claim.

\FloatBarrier

\section{Protection Sensitivity and Boundary Checks}
\label{app:protection-supp}
\label{app:protection-details}

\subsection{Common repair harness}
\label{app:guard-proof}
\label{app:repair-protocol}

Every active arm begins from the same first criterion-reaching parent
checkpoint. It selects exactly $K_p$ eligible discovery examples, but shares
the neutral stream $\mathcal R_0$, example order, optimizer random stream,
processed-token budget, projection, guards, and save schedule. The neutral
stream enters $L_{\rm rep}$ through the entry-anchored KL term in
Eq.~\eqref{eq:budgeted-repair}; no susceptibility score selects it. The token
budget charges protected, neutral, constraint-refresh, guard, and
checkpoint-evaluation tokens.

Projection covers only differentiable surrogate constraints. Each is oriented
as $c_j(\theta)\leq0$ on a frozen guard set and is active when
$c_j(\theta)\geq-\kappa_j$. The matrix $G$ is refreshed every $m_{\rm ref}$
accepted steps and Eq.~\eqref{eq:repair-projection} uses ridge $\lambda$.
For a forget token $(i,s)$, the anchored adverse margin is
\[
 r^{f}_{is}(\theta)=\log p_\theta(y_{is}\mid y_{i,<s})-
 \log p_{t^\dagger}(y_{is}\mid y_{i,<s});
\]
for a neutral or utility token it is the NLL increase
\[
 r^{u}_{is}(\theta)=-\log p_\theta(y_{is}\mid y_{i,<s})+
 \log p_{t^\dagger}(y_{is}\mid y_{i,<s}).
\]
A proposed step must keep every token margin below its frozen
$\epsilon^{\rm tok}$ and every example-average margin below
$\epsilon^{\rm ex}$. Otherwise the parameters are restored, the step size is
halved, and the step is retried up to $J_{\rm retry}$ times; exhaustion stops
the arm. The manifest freezes $\beta$, $\kappa_j$, $m_{\rm ref}$, $\lambda$,
both tolerances, $J_{\rm retry}$, and guard IDs before target runs.

Non-differentiable extraction, generation, and full utility metrics are not
represented in $G$. They are evaluated only at the shared save schedule on
selector-safe sets. The reported output is the last checkpoint satisfying
every row of Table~\ref{tab:constraint-contract}; if none does, the arm is
infeasible. Sealed $\Delta$NLL and native retained outcomes are opened only
after this selection. The artifact records refreshes, retries, accepted and
rejected updates, margins, and the selected step.

Repeated random allocation uses $B_{\rm rand}$ pools frozen before outcomes.
Every draw must complete and satisfy the same constraints. Draw IDs are
resampled and averaged inside their trajectory within each bootstrap
replicate. Diagnostic arms replace only the selector: exact-energy allocation
substitutes $q_G^\star$ for $\widehat q_G$; prediction-weight allocation uses
$\widehat\alpha_{\rm pred}$; uniform rehearsal ignores ranking; and a
realized-discovery-damage oracle gives an unattainable repair-capacity
reference. All remain outside the four-comparator claim.

\subsection{Budget and specificity checks}
\label{app:specificity}

The specificity analysis keeps the sealed target profile fixed. Other-request
motion uses a size-matched forget set; random motion matches the target block
displacement and uses directions independent of profiling. These controls
bound interpretation: loss-shake energy is intentionally direction- and
request-agnostic, whereas proximity should carry the request-specific part.
MUSE-Books and PISTOL remain boundary settings and never count toward the
paper-level breadth rule.

Hyperparameter grids, source-to-target weight carryover, exploratory traces of
directional persistence, and early hypothesis development are artifact
diagnostics rather than paper tables. They never select a target rule or fill
a confirmatory cell.

\begin{table*}[!t]
\caption{\textbf{Protection sensitivity and specificity boundaries.} Panel A
varies the protected fraction under the unchanged five-arm contract. Panel B
asks whether each frozen component responds specifically to the target
request. Primary parent decisions and their decisive worst UCB remain in
Table~\ref{tab:core-evidence}; complete arm effects, constraint vectors, guard
events, and random-pool rows are released with the artifact.}
\label{tab:crossed}
\label{tab:budget-sweep}
\label{tab:protection-ablations}
\label{app:budget-sweep}
\centering
\scriptsize
\setlength{\tabcolsep}{1.8pt}
\resizebox{0.72\textwidth}{!}{%
\begin{tabular}{@{}lccccc@{}}
\toprule
\multicolumn{6}{@{}l}{\textbf{A. Fixed-budget sensitivity}} \\
Budget & Eligible/pass & worst $\Delta$ [UCB] & Bottleneck &
min F/U slack & Updates/support \\
\midrule
0.05 & \tblph & \tblph & \tblph & \tblph/\tblph & \tblph/\tblph \\
0.10 (primary) & \tblph & \tblph & \tblph & \tblph/\tblph & \tblph/\tblph \\
0.20 & \tblph & \tblph & \tblph & \tblph/\tblph & \tblph/\tblph \\
\bottomrule
\end{tabular}}

\vspace{3pt}
\resizebox{0.72\textwidth}{!}{%
\begin{tabular}{@{}llcccc@{}}
\toprule
\multicolumn{6}{@{}l}{\textbf{B. Request specificity and motion controls}} \\
Motion & Conditioning & $\Delta\rho_{G/H/S}$ & top-$q$ lift & Disp. matched & Support \\
\midrule
Target unlearning & target request & \tblph/\tblph/\tblph & \tblph & yes & \tblph \\
Other-request & different request & \tblph/\tblph/\tblph & \tblph & yes & \tblph \\
Benign SFT & neutral stream & \tblph/\tblph/\tblph & \tblph & yes & \tblph \\
Random motion & independent directions & \tblph/\tblph/\tblph & \tblph & yes & \tblph \\
\bottomrule
\end{tabular}}
\label{tab:specificity}
\end{table*}

\FloatBarrier
